\section*{Hoofdstuk \ref{ch:g11:becoming-male-female} --- Man of vrouw worden}

\begin{solution}{exo:g11:becoming-male-female:1}
Die van de vader: elke eicel draagt een X, terwijl de helft van de
zaadcellen een X en de helft een Y draagt, zodat XX en XY even
waarschijnlijk zijn.
\end{solution}

\begin{solution}{exo:g11:becoming-male-female:2}
Twee onbepaalde geslachtsklieren, elk met twee gangen ernaast: een gang
van Wolff (de toekomstige mannelijke geslachtsweg) en een gang van Müller
(de toekomstige vrouwelijke). Gelijk bij XX- en XY-embryo's.
\end{solution}

\begin{solution}{exo:g11:becoming-male-female:3}
Op het \omterm{def:g11:cell-cycle-mitosis:chromosome}{Y-chromosoom}; het komt rond week 7 in de onbepaalde
geslachtsklier tot uiting, en zijn \omterm{def:g11:gene-expression:protein}{eiwit} schakelt de \omterm{def:g10:universal-dna:gene}{genen} aan die van de
geslachtsklier een teelbal maken.
\end{solution}

\begin{solution}{exo:g11:becoming-male-female:4}
\omterm{prop:g11:becoming-male-female:hormones}{Testosteron} houdt de gangen van Wolff in stand en laat ze uitgroeien tot
de mannelijke geslachtsweg; anti-Müllers hormoon laat de gangen van
Müller wegschrompelen.
\end{solution}

\begin{solution}{exo:g11:becoming-male-female:5}
Een centrum in de hersenen geeft stoten af van een hormoon dat de
hypofyse LH en FSH laat afscheiden, die de geslachtsklieren aanzetten.
Secundaire geslachtskenmerken: de veranderingen die de geslachtshormonen
voortbrengen --- lichaamsbeharing, de stem, de borsten, de verdeling van
spier en vet, de groeispurt, de groei van de voortplantingsorganen.
\end{solution}

\begin{solution}{exo:g11:becoming-male-female:6}
Hij ontwikkelde een vrouwelijke geslachtsweg, ondanks zijn \omterm{def:g11:cell-cycle-mitosis:chromosome}{Y-chromosoom}:
de mannelijke geslachtsweg vergt de hormonen van de teelbal, en de
vrouwelijke ontstaat zonder enige geslachtsklier.
\end{solution}

\begin{solution}{exo:g11:becoming-male-female:7}
\omterm{prop:g11:becoming-male-female:hormones}{Testosteron} bouwt de mannelijke gangen maar ruimt de vrouwelijke niet op;
daarvoor moet de teelbal een tweede stof afscheiden --- AMH.
\end{solution}

\begin{solution}{exo:g11:becoming-male-female:8}
Teelballen (Sry werkt); gangen van Wolff behouden, die van Müller
weggeschrompeld; mannelijke uitwendige bouw. Onvruchtbaar: het maken van
zaadcellen vergt \omterm{def:g10:universal-dna:gene}{genen} van de Y die het transgen niet levert.
\end{solution}

\begin{solution}{exo:g11:becoming-male-female:9}
Oestrogenen: 50\% rond 11, plateau rond 14--15. \omterm{prop:g11:becoming-male-female:hormones}{Testosteron}: 50\% rond 13,
plateau rond 16--17.
\end{solution}

\begin{solution}{exo:g11:becoming-male-female:10}
De botten houden op met groeien wanneer het geslachtshormoon zijn
volwassen niveau bereikt; oestrogenen bereiken dat ongeveer twee jaar
eerder dan testosteron, dus sluiten de groeischijven van meisjes eerder.
\end{solution}

\begin{solution}{exo:g11:becoming-male-female:11}
Geen teelbal: de geslachtsklier wordt een eierstokachtige klier; geen
testosteron of AMH: de gangen van Müller ontwikkelen zich, die van Wolff
schrompelen weg, en de uitwendige bouw is bij de geboorte vrouwelijk. In
de \omterm{prop:g11:becoming-male-female:puberty}{puberteit} scheiden de geslachtsklieren, die geen normale eicellen
bevatten, weinig af; zonder hormoonbehandeling ontwikkelen de secundaire
kenmerken zich slecht.
\end{solution}

\begin{solution}{exo:g11:becoming-male-female:12}
Teelballen (\omterm{prop:g11:becoming-male-female:sry}{SRY} werkt). De gangen van Wolff schrompelen weg (geen reactie
op testosteron); de gangen van Müller schrompelen weg (AMH werkt): geen
inwendige geslachtsweg van beide typen. De uitwendige bouw is
vrouwelijk (haar mannelijke vorm vergt de werking van testosteron). In de
\omterm{prop:g11:becoming-male-female:puberty}{puberteit} brengt testosteron dat in oestrogenen wordt omgezet borsten en
een vrouwelijke bouw voort, zonder lichaamsbeharing.
\end{solution}

\begin{solution}{exo:g11:becoming-male-female:13}
Zodra een geslachtsklier een teelbal of een eierstok is geworden, kunnen
haar hormonen de gangen en het lichaam net zo vormen als bij zoogdieren;
wat verschilt is de eerste stap --- een temperatuurgevoelig proces
vervangt \omterm{prop:g11:becoming-male-female:sry}{SRY} als de schakelaar die over de geslachtsklier beslist.
\end{solution}

\begin{solution}{exo:g11:becoming-male-female:14}
De X draagt honderden \omterm{def:g10:universal-dna:gene}{genen} die elke \omterm{def:g10:cells-common-unit:cell}{cel} nodig heeft, dus overleeft geen
enkel embryo zonder er één; de Y draagt vrijwel niets behalve de \omterm{def:g10:universal-dna:gene}{genen}
voor het bepalen van de teelbal en voor de zaadcellen, dus missen
XX-personen niets wat zij nodig hebben. Een Y gaat alleen naar zonen
omdat dochters de X van hun vader krijgen.
\end{solution}

\begin{solution}{exo:g11:becoming-male-female:15}
Het betekent dat de vrouwelijke geslachtsweg zich ontwikkelt wanneer geen
hormoon van de teelbal tussenbeide komt --- een feit over welke signalen
nodig zijn, geen rangorde. Het weerspiegelt een gemeenschappelijk
embryonaal plan waarop één cellijn, de teelbal, met twee hormonen een
alternatief oplegt: de twee geslachten zijn twee uitkomsten van één
programma, die door een schakelaar verschillen.
\end{solution}

\begin{solution}{pb:g11:becoming-male-female:1}
\textbf{1.} De eierstok is niet nodig voor de vrouwelijke geslachtsweg,
die zonder enige geslachtsklier ontstaat; de teelbal is nodig voor de
mannelijke geslachtsweg en voor het opruimen van de vrouwelijke gangen.

\textbf{2.} Na dag 22 zijn de gangen al onder invloed van de
geslachtsklier begonnen te differentiëren; de proef moet vóór de
beslissing ingrijpen.

\textbf{3.} De hormonen van de teelbal, door het bloed naar de gangen
gebracht: het ent maskuliniseert ze ook van veraf.

\textbf{4.} \omterm{prop:g11:becoming-male-female:hormones}{Testosteron} bouwt de mannelijke gangen; het ruimt de
vrouwelijke niet op.

\textbf{5.} Een tweede stof van de teelbal moet de gangen van Müller
laten wegschrompelen; dat was te bevestigen door haar uit teelballen te
winnen en alleen in te brengen, met de verwachting dat de vrouwelijke
gangen verdwijnen en de mannelijke niet ontstaan --- en dat is wat AMH
doet.

\textbf{6.} Aan de geopereerde kant, zonder testosteron of AMH: een
vrouwelijke gang, geen mannelijke. Aan de ongeschonden kant: een
mannelijke gang, de vrouwelijke weggeschrompeld.

\textbf{7.} Aan de kant van het ent een mannelijke gang en geen
vrouwelijke; aan de andere kant alleen een vrouwelijke gang.

\textbf{8.} Aan de kant van het kristal beide gangen aanwezig; aan de
andere kant alleen de vrouwelijke gang.

\textbf{9.} Elke foetus dient als zijn eigen controle: het verschil
tussen de twee kanten kan alleen komen door de plaatselijke aanwezigheid
van de teelbal of het testosteron, of door hun afwezigheid.

\textbf{10.} Haar eierstokken zijn verwijderd of overstemd, en het geënte
teelbaltje, van een ander dier, maakt in een vreemde gastheer geen
zaadcellen; de geslachtsweg is mannelijk maar er worden geen gameten
gemaakt.

\textbf{11.} \omterm{prop:g11:becoming-male-female:sry}{SRY} aanwezig: een teelbal; testosteron en AMH: mannelijke
geslachtsweg, mannelijke bouw; in de \omterm{prop:g11:becoming-male-female:puberty}{puberteit} testosteron: mannelijke
kenmerken. Onvruchtbaar (geen \omterm{def:g10:universal-dna:gene}{genen} van de Y voor zaadcellen).

\textbf{12.} Geen werkend \omterm{prop:g11:becoming-male-female:sry}{SRY}: geen teelbal; geen hormonen: vrouwelijke
geslachtsweg en bouw; slecht ontwikkelde geslachtsklieren.

\textbf{13.} Teelballen; gangen van Wolff ontwikkeld (testosteron werkt);
gangen van Müller blijven ook bestaan (geen reactie op AMH): een man met
daarbij een baarmoeder en eileiders; mannelijke bouw en \omterm{prop:g11:becoming-male-female:puberty}{puberteit}.

\textbf{14.} Alleen het derde: werkende teelballen en mannelijke gangen,
al kan de blijvende baarmoeder het indalen van de teelballen bemoeilijken.
Het eerste mist de \omterm{def:g10:universal-dna:gene}{genen} voor zaadcellen; het tweede heeft geen werkende
geslachtsklieren.

\textbf{15.} Het chromosomale geslacht voorspelt het anatomische alleen
via de stappen ertussen: een werkend \omterm{prop:g11:becoming-male-female:sry}{SRY}, werkende hormonen, werkende
receptoren. Verander welke stap dan ook en de twee niveaus lopen uiteen.

\textbf{16.} Meisjes: 20\% op 10, spurt vanaf ongeveer 11, stop rond 15.
Jongens: 20\% rond 12.5, spurt vanaf ongeveer 13.5, stop rond 17.

\textbf{17.} Meisjes ongeveer 4 jaar, jongens ongeveer 3.5 --- maar
jongens beginnen vanaf een grotere lengte, doordat zij twee jaar langer
met het kindertempo zijn doorgegroeid.

\textbf{18.} Jongens ongeveer $3.5 \times 8 = \qty{28}{cm}$, meisjes $4 \times
7 = \qty{28}{cm}$: de spurts zijn vergelijkbaar; het verschil
bij volwassenen komt vooral van de twee extra jaren kindergroei vóór de
spurt van de jongens.

\textbf{19.} Een vroege stijging brengt de groeischijven jaren te vroeg
tot sluiting: het kind groeit nu snel maar stopt op een geringe volwassen
lengte. De hormonen uitstellen verlengt de groei van de kindertijd.

\textbf{20.} \omterm{prop:g11:becoming-male-female:hormones}{Testosteron}, dat de mannelijke gangen bouwt, en
anti-Müllers hormoon, dat de vrouwelijke opruimt; beide gemaakt door de
teelbal, een orgaan dat door het ene \omterm{def:g10:universal-dna:gene}{gen} \omterm{prop:g11:becoming-male-female:sry}{SRY} wordt aangezet.
\end{solution}
